\documentclass{article}
\usepackage[utf8]{inputenc}
\usepackage{amsmath}
\usepackage{geometry}
\geometry{margin=2cm}
\begin{document}

\section*{Questão 167 - ENEM 2024 - Matemática}

\subsection*{Enunciado}
Atualmente, há telefones celulares com telas de diversos tamanhos e em formatos retangulares. Alguns deles apresentam telas medindo \(3 \tfrac{1}{2}\) polegadas, com determinadas especificações técnicas. Além disso, em muitos modelos, com a inclusão de novas funções no celular, suas telas ficaram maiores, sendo comum encontrarmos atualmente telas medindo \(4 \tfrac{5}{6}\) polegadas, conforme a figura.

\begin{center}
\textit{Disponível em: www.tecmundo.com.br. Acesso em: 5 nov. 2014 (adaptado).}
\end{center}

A diferença de tamanho, em valor absoluto, entre as medidas, em polegada, das telas do celular 2 e do celular 1, representada apenas com uma casa decimal, é:

\bigskip

\includegraphics[width=0.6\textwidth]{imagem_celulares_enunciado.png}

\subsection*{Solução Técnica}

A resposta correta é \textbf{1,3} polegadas, que corresponde à alternativa D.

\textbf{Estratégia:} Converter as medidas de frações mistas para números decimais, calcular a diferença absoluta e arredondar para uma casa decimal.

\textbf{Passos:}
\begin{enumerate}
  \item \(3 \tfrac{1}{2} = 3 + 0,5 = 3,5\) polegadas
  \item \(4 \tfrac{5}{6} = 4 + 0,8333 = 4,8333\) polegadas
  \item Diferença: \(|4,8333 - 3,5| = 1,3333\)
  \item Arredondar: \(1,3\) polegadas
\end{enumerate}

\subsection*{Fórmula Utilizada}

\[
\text{Diferença absoluta} = \left| \text{medida do celular 2} - \text{medida do celular 1} \right|
\]

\subsection*{Explicação Didática}
Para resolver a questão, convertemos as frações mistas para decimais, depois calculamos a diferença absoluta entre as medidas e arredondamos adequadamente.

\subsection*{Análise dos Distratores}
As alternativas incorretas se baseiam em erros comuns como:
\begin{itemize}
  \item Confusão na conversão de frações para decimais (calculando diferenças muito pequenas).
  \item Ignorar o valor da fração \(\frac{5}{6}\).
  \item Arredondamento incorreto.
  \item Erro na operação, como somar ao invés de subtrair.
\end{itemize}

\subsection*{Perfil da Questão}
A questão trabalha habilidades de conversão de fração mista em decimal, cálculo de diferença absoluta e arredondamento, cobrando raciocínio matemático de nível médio com interpretação visual.

\bigskip

\textbf{Confiança da análise:} Alta

\end{document}